\subsection{例}

\begin{enumerate}
\def\labelenumi{\arabic{enumi})}
\item
  设 E 是一个赋范空间。给 E 带上由 ab = 0（对 E 中所有的 a 与 b）定义的乘法；这样就在 E 上定义了一个赋范代数结构。
\item
  设 X 是一个集合。我们用 \(\mathcal{B}(X;K)\) 表示 X 上取值于 K 的有界函数所成的赋范代数，带上范数

\[
\| f \| = \sup _ {x \in X} | f (x) |
\]

(TG, X, p.~22)。它是 K 上的一个幺 Banach 代数 (TG, X, p.~21, 推论 1)。它是交换的。设 \(f\) 是 \(\mathcal{B}(\mathrm{X};\mathrm{K})\) 的一个元素。则 \(f\) 在 \(\mathcal{B}(\mathrm{X};\mathrm{K})\) 中可逆必须且只需

\[
\inf _ {x \in X} | f (x) | > 0.
\]

于是 \(f\) 的谱是由使下式成立的 \(\lambda\in K\) 所成的集合：

\[
\inf _ {x \in \mathrm{X}} | f (x) - \lambda | = 0,
\]

也就是值集 \(f(\mathrm{X})\)（\(f\) 的值集）在 K 中的闭包。

\item
  设 X 是一个拓扑空间。我们用 \(\mathcal{C}_{b}(\mathrm{X};\mathrm{K})\) 表示 \(\mathcal{B}(\mathrm{X};\mathrm{K})\) 的由 X 上取值于 K 的连续且有界的函数所成的幺子代数，用 \(\mathcal{C}_{0}(\mathrm{X};\mathrm{K})\) 表示 \(\mathcal{C}_{b}(\mathrm{X};\mathrm{K})\) 的由在无穷远处趋于 0 的函数所成的子代数 (参见 INT, III, §1, 第 2 小节及 TG, X, p.~40, 推论 2)。我们回顾，\(\mathcal{H}(\mathrm{X};\mathrm{K})\) 表示 \(\mathcal{C}_{b}(\mathrm{X};\mathrm{K})\) 的由具有紧支集的连续函数所成的子代数。
代数 \(\mathcal{C}_{b}(\mathrm{X};\mathrm{K})\) 与 \(\mathcal{C}_{0}(\mathrm{X};\mathrm{K})\) 是 K 上的交换 Banach 代数；事实上，它们是 \(\mathcal{B}(\mathrm{X};\mathrm{K})\) 的闭子空间 (TG, X, p.~21, 推论 2 及 INT, III, §1, 第 2 小节)。

由于一个处处非零的连续函数的逆是连续的，子代数 \(\mathcal{C}_{b}(\mathrm{X};\mathrm{K})\) 是 \(\mathcal{B}(\mathrm{X};\mathrm{K})\) 的一个全子代数。特别地，\(\mathcal{C}_{b}(\mathrm{X};\mathrm{K})\) 的元素 f 的谱等于 \(\overline{f(\mathrm{X})}\)。

若 X 是离散的，则有 \(\mathcal{C}_b(\mathrm{X};\mathrm{K}) = \mathcal{B}(\mathrm{X};\mathrm{K})\)。

若 X 是紧的，则有 \(\mathcal{C}_{b}(\mathrm{X};\mathrm{K}) = \mathcal{H}(\mathrm{X};\mathrm{K}) = \mathcal{C}(\mathrm{X};\mathrm{K})\)，即 X 上取值于 K 的连续函数所成的赋范幺代数 \(\mathcal{C}(\mathrm{X};\mathrm{K})\)。这时，元素 \(f \in \mathcal{C}(\mathrm{X};\mathrm{K})\) 可逆必须且只需 f 不取值 0，且 f 的谱等于 f 的值集 \(f(\mathrm{X})\)。

此后假定 X 不是紧的。于是代数 \(\mathcal{C}_{0}(X;K)\) 不是幺的；由它添加单位元而得到的代数与子代数 \(K \cdot 1 \oplus \mathcal{C}_{0}(X;K)\)（\(\mathcal{C}(X;K)\) 的由在无穷远处有极限的函数所成的子代数）视为同一。这个子代数的一个元素可逆必须且只需它不为零且它在无穷远处的极限非零。由此得出，对每个 \(f \in \mathcal{C}_{0}(X;K)\)，f 的谱等于 \(f(X) \cup \{0\}\)。

\(f \in \mathcal{K}(\mathrm{X};\mathrm{K})\) 的谱等于 \(f(\mathrm{X})\)（事实上，若 \(\mathrm{X}\) 不是紧的，则 0 属于 \(f(\mathrm{X})\)）。

\item
  设 \(n \geqslant 0\) 是一个整数。设 \(A_{n}\) 是由函数 \(f: [0,1] \to K\)（在 \([0,1]\) 上有直到 \(n\) 阶的连续导数）所成的代数，带上范数
\[
\| f \| = \sum_ {k = 0} ^ {n} \frac {1}{k !} \sup _ {0 \leqslant t \leqslant 1} | f ^ {(k)} (t) |.
\]

若 \(f,g\in \mathrm{A}_n\)，则有

\[
\begin{array}{l} \| f g \| = \sum_ {k = 0} ^ {n} \frac {1}{k !} \sup | (f g) ^ {(k)} (t) | = \sum_ {k = 0} ^ {n} \frac {1}{k !} \sup \Bigl | \sum_ {s = 0} ^ {k} \binom {k} {s} f ^ {(s)} (t) g ^ {(k - s)} (t) \Bigr | \\ \leqslant \sum_ {k = 0} ^ {n} \sum_ {s = 0} ^ {k} \frac {1}{s ! (k - s) !} \sup _ {0 \leqslant t \leqslant 1} | f ^ {(s)} (t) | \sup _ {0 \leqslant t \leqslant 1} | g ^ {(k - s)} (t) | = \| f \|   \| g \|, \end{array}
\]

这是由 Leibniz 公式 (FVR, I, p.~28, 命题 2) 得出的，因而 \(A_{n}\) 是一个交换的幺 Banach 代数。

\item
  设 E 是一个 Banach 空间，其范数记作 p。E 的连续自同态所成的代数 \(\mathcal{L}(\mathrm{E})\)，带上范数
\[
\| u \| = \sup _ {p (x) \leqslant 1} p (u (x))
\]

是一个幺 Banach 代数 (EVT, III, p.~14 及 p.~24, 推论 2; TG, X, p.~23, 公式 (3))。

\item
  设 G 是一个局部紧群。用 e 表示它的单位元。设 \(\mathcal{M}^{1}(G)\) 是 G 上的有界复测度所成的 Banach 空间 (INT, III, p.~57)。卷积积 (INT, VIII, p.~120, 定义 1) 给 \(\mathcal{M}^{1}(G)\) 带上一个复 Banach 代数结构 (INT, VIII, §3, 第 1 小节, 命题 2)，其单位元是由置于点 e 的单位质量所定义的测度 \(\varepsilon_{e}\)。若 G 交换，这个 Banach 代数是交换的。具有紧支集的测度所成的空间 \(\mathcal{C}'(G)\) 是 \(\mathcal{M}^{1}(G)\) 的一个子代数 (同前)。
\item
  设 G 是一个带有 Haar 测度 \(\mu\) 的局部紧群。则 \(\mathrm{L}_{\mathrm{K}}^{1}(\mathrm{G},\mu)\) 关于卷积积是一个 Banach 代数 (INT, VIII, 命题 12, p.~166)。若 K = C，则由 \(f \mapsto f \mu\) 定义的映射使我们可以把 \(\mathrm{L}_{\mathbf{C}}^{1}(\mathrm{G},\mu)\) 与 Banach 代数 \(\mathcal{M}^{1}(\mathrm{G})\) 的一个子代数视为同一。
若 G 交换，则 Banach 代数 \(\mathrm{L}_{\mathrm{K}}^{1}(\mathrm{G},\mu)\) 是交换的。

\item
  在例 7 中取 G = Z 及 K = C。则 \(\mathrm{L}_{\mathbf{C}}^{1}(\mathbf{Z})\) 是由序列 \((x_{n})_{n\in\mathbf{Z}}\)（满足 \(\sum_{n}|x_{n}|<+\infty\)）所成的复交换 Banach 代数，元素 \((x_{n})\) 与 \((y_{n})\) 的乘积是 \((z_{n})\)，其中
\[
z _ {n} = \sum_ {k \in {\bf Z}} x _ {k} y _ {n - k}
\]

而范数为 \(\|(x_{n})\|=\sum_{n}|x_{n}|\)。这个代数以序列 \(\varepsilon=(\varepsilon_{n})\) 为单位元，其中 \(\varepsilon_{0}=1\)，且 \(\varepsilon_{n}=0\)（对 \(n\neq0\)）。

用 U 表示 C 中的单位圆。若 \(x = (c_{n})\) 是 A 的一个元素，用 \(\varphi(x)\) 表示 U 上的连续函数，它在 \(e^{it}\) 处的值为

\[
\varphi (x) (e ^ {i t}) = \sum_ {n \in {\bf Z}} c _ {n} e ^ {i n t}.
\]

可以验证，\(\varphi\) 是从 \(\mathrm{L}_{\mathbf{C}}^{1}(\mathbf{Z})\) 到 \(\mathbf{U}\) 上连续函数的一个代数 A 上的态射，A 中的乘法就是通常的乘法。把等式逐项积分

\[
\left(\sum_ {m \in {\bf Z}} c _ {m} e ^ {i m t}\right) \cdot e ^ {- i n t} = \varphi (x) (e ^ {i t}) \cdot e ^ {- i n t},
\]

就得到

\[
c _ {n} = \frac {1}{2 \pi} \int_ {0} ^ {1} \varphi (x) (e ^ {i t}) e ^ {- i n t} d t.
\]

特别地，由此得出态射 \(\varphi\) 是单射。代数 A 带上由 \(\mathrm{L}_{\mathbf{C}}^{1}(\mathbf{Z})\) 的范数经由 \(\varphi\) 诱导的范数，称为绝对收敛 Fourier 级数的 Banach 代数。它以函数 \(1 = \varphi(\varepsilon)\) 为单位元。

\item
  设 \(\Delta\) 是由复数 \(z\)（满足 \(|z| \leqslant 1\)）所成的圆盘。\(\Delta\) 上连续且在 \(\Delta\) 的内部全纯的函数所成的代数 (VAR, R1, p.~26, 3.2.1) 带上范数 \(\|f\| = \sup_{z \in \Delta} |f(z)|\)。则 A 是一个交换的幺 Banach 代数。
\end{enumerate}

